\pdfminorversion=7
\documentclass[12pt,letterpaper]{article}

\usepackage[T1]{fontenc}
\usepackage[utf8]{inputenc}
\usepackage[american]{babel}
\usepackage{csquotes}

\usepackage[
    style=science,
    backend=biber,
    sorting=none,
    articletitle=true,
    url=false,
    doi=false,
    eprint=false,
    isbn=false
]{biblatex}
\addbibresource{references.bib}
\AtEveryBibitem{\clearlist{language}}

\usepackage{newtxtext,newtxmath}
\usepackage{microtype}
\usepackage{amsmath}
\usepackage{graphicx}
\usepackage{booktabs}
\usepackage[margin=0.92in,top=0.82in,bottom=0.95in,headsep=18pt,footskip=24pt]{geometry}
\usepackage[font=small,labelfont=bf,labelsep=period]{caption}
\usepackage{subcaption}
\usepackage{float}
\usepackage{afterpage}
\usepackage{array}
\usepackage{multirow}
\usepackage{setspace}
\usepackage{titlesec}
\usepackage{enumitem}
\usepackage{fancyhdr}
\usepackage{needspace}
\usepackage{etoolbox}
\usepackage[section]{placeins}
\usepackage{flafter}
\usepackage[table]{xcolor}

\definecolor{ClayLine}{HTML}{EEDBDD}
\definecolor{SumiInk}{HTML}{2D2426}
\definecolor{DustyTaupe}{HTML}{866A6D}

\usepackage[hidelinks]{hyperref}
\usepackage[switch]{lineno}

\color{SumiInk}
\setstretch{1.45}
\setlength{\parindent}{0pt}
\setlength{\parskip}{0.42em}
\setlength{\textfloatsep}{8pt plus 1pt minus 1pt}
\setlength{\floatsep}{7pt plus 1pt minus 1pt}
\setlength{\intextsep}{7pt plus 1pt minus 1pt}
\setlength{\abovecaptionskip}{5pt}
\setlength{\belowcaptionskip}{2pt}
\setlength{\tabcolsep}{9pt}
\renewcommand{\arraystretch}{1.14}
\captionsetup{justification=raggedright,singlelinecheck=false}
\setlist[itemize]{leftmargin=1.3em,itemsep=0.18em,topsep=0.25em}
\setcounter{topnumber}{3}
\setcounter{bottomnumber}{2}
\setcounter{totalnumber}{4}
\renewcommand{\topfraction}{0.92}
\renewcommand{\bottomfraction}{0.85}
\renewcommand{\textfraction}{0.08}
\renewcommand{\floatpagefraction}{0.8}
\setlength{\headheight}{24pt}
\emergencystretch=2em
\raggedbottom

\titleformat{\section}
  {\Large\sffamily\bfseries\color{SumiInk}}
  {\thesection}
  {0.6em}
  {}
  [\vspace{0.25em}\color{ClayLine}\titlerule]
\titleformat{\subsection}
  {\large\sffamily\bfseries\color{SumiInk}}
  {\thesubsection}
  {0.55em}
  {}
\titleformat{\subsubsection}
  {\normalsize\sffamily\bfseries\color{SumiInk}}
  {}
  {0em}
  {}
\titlespacing*{\section}{0pt}{1.3ex plus 0.4ex minus 0.2ex}{0.7ex}
\titlespacing*{\subsection}{0pt}{1.15ex plus 0.2ex minus 0.2ex}{0.45ex}
\titlespacing*{\subsubsection}{0pt}{0.9ex plus 0.15ex minus 0.1ex}{0.25ex}

\pretocmd{\subsubsection}{\needspace{4\baselineskip}}{}{}
\pretocmd{\subsection}{\needspace{5\baselineskip}}{}{}

\pagestyle{fancy}
\fancyhf{}
\fancyhead[L]{\footnotesize\sffamily\color{SumiInk} Low-threat game learning}
\fancyhead[R]{\footnotesize\sffamily\color{DustyTaupe}\nouppercase{\leftmark}}
\fancyfoot[C]{\footnotesize\sffamily\thepage}
\renewcommand{\headrulewidth}{0.35pt}
\renewcommand{\headrule}{\hbox to\headwidth{\color{ClayLine}\leaders\hrule height \headrulewidth\hfill}}
\renewcommand{\sectionmark}[1]{\markboth{#1}{}}

\makeatletter
\def\fps@figure{tbp}
\def\fps@table{tbp}
\makeatother
\newcommand{\backmattersection}[1]{\section*{#1}\markboth{#1}{}}

\begin{document}
\linenumbers
\thispagestyle{empty}

\begin{center}
    \vspace*{-0.6em}

    {\fontsize{20}{22}\selectfont\bfseries\color{SumiInk}
    A low-threat practice-room interface for neurodivergent game learning\par}

    \vspace{0.95em}

    {\normalsize\color{SumiInk}
    Alan N. Pham$^{1,*}$\par}

    \vspace{0.38em}

    {\small\color{DustyTaupe}
    $^{1}$AO Labs\\
    $^{*}$Correspondence: \href{https://aolabs.io}{aolabs.io}\par}
\end{center}

\vspace{0.35em}
\noindent\textbf{Fast multiplayer games can convert practice into threat when social judgment, sensory load, ambiguous damage, and repeated failure are experienced as one continuous demand. This manuscript reports a single-user design study of \href{https://league.aolabs.io}{league.aolabs.io}, a public review room for League of Legends practice. The system does not claim clinical efficacy, population generality, a guaranteed rank outcome, or competitive optimization. It separates two kinds of Samira evidence that serve different jobs. A public screenshot library preserves exact tip images and derives source-grounded tips for later review and morning-email rotation. A separate coach-game-analysis intake stores the coach's complete response, extracts game facts and a structured VOD entry, and remains the only source for note rank reads, rank and CS@10 charts, and note PDFs. One visible \emph{copy coach message} button copies the complete review request that the user sends with a video through an external coach channel; the website does not upload or send that video. The strongest contribution is a design pattern rather than a gameplay result: keep capture dense and direct, preserve the primary source before generated analysis, and keep advice sources separate enough that a screenshot can improve the morning cue without changing the performance record. This bounded case suggests how personal learning software can support high-repetition play without becoming another source of evaluation pressure.}

\newpage
\pagestyle{fancy}

\section{Introduction}

Expert performance depends on repeated, structured practice, but practice is only useful when the learner can remain in the task long enough for repetitions to accumulate \cite{ericsson1993}. In competitive games, especially player-versus-player games, the immediate source of failure is not only mechanical difficulty. The learner must also manage uncertainty, sensory load, social interpretation, and the felt cost of visible mistakes. Anxiety can impair attentional control \cite{eysenck2007}; cognitive load can interfere with learning when too many elements must be processed at once \cite{sweller1988}. A player who is trying to learn a champion, lane, camera, item build, opponent behavior, teammate expectations, and self-regulation at the same time may experience the game as a judgment environment rather than a practice environment.

The system described here began from a simple transfer from music practice: keep the task playable, choose one thing, trust repetitions, and let improvement compound. The same principle was applied to League of Legends. The first rule was not to win lane or outplay an opponent; it was to farm and survive. That rule later compressed into the phrase ``next safe minion.'' As the practice record grew, the problem became less about collecting more advice and more about making advice reviewable without overwhelming the player before queueing.

\href{https://league.aolabs.io}{league.aolabs.io} is a public AO Labs page built for that purpose. It is not a Riot Games product, a general League guide, or an official ranked-MMR system. The current home page first exposes a Samira tip-image intake and then a distinct coach-game-analysis intake, notes-only rank and CS@10 charts, and dense saved-game entries. Earlier Samira, Caitlyn, and Fizz review-room material remains source history. The design is source-bounded: it uses exact user-supplied screenshots, the user's supplied chat arc, public coach responses, local recordings, local client reads, and saved Samira note blocks as its corpus, and it avoids claims that would require clinical, population-scale, or competitive outcome evidence. Challenger is a direction for skill development, not a prediction or promise.

\section{Results}

\subsection{A long learning arc was reduced to two direct capture paths}

The seed corpus contained a long narrative account of League practice. The arc included lane fear, camera stability, sensory control, champion feel, bot practice, death tolerance, public yapping as pressure regulation, and champion-specific rules for Samira, Caitlyn, and Fizz. The current design result is not a long visible essay, a motivational hero, or a broad champion library. It is two source-specific capture paths. The first receives screenshots containing general Samira advice, preserves the exact image, and derives grounded tips that can rotate into the morning email. The second prepares an external VOD-review request and receives the coach's completed analysis as a game record. This second path alone supplies the parsed date, queue, final Alan/Samira facts, explicit gameplay-estimated rank, notes-only rank trend, separate CS@10 trend, and entry PDF. CS@10 remains early-farm evidence rather than silently changing the rank score, and a screenshot can never become performance evidence.

This separation is the main interface result. It lets general guidance improve the next morning cue without contaminating the history of analyzed games. It also makes the page compatible with the user's stated constraint that the surface be beautiful, dense, and immediately usable. The long record remains useful as paper, log, and archive context, but the first screen does not ask the user to process a giant review tray, generic champion taxonomy, or large empty analytics frame before adding a source.

\begin{table}[H]
\centering
\caption{Compression of the source arc into champion-specific review units.}
{\small
\setlength{\tabcolsep}{5pt}
\begin{tabular}{p{0.24\linewidth}p{0.30\linewidth}p{0.30\linewidth}}
\toprule
Source burden & Review unit & Interface expression \\
\midrule
Enemy lane feels like judgment & Archived situation & Next safe minion remains in the paper record \\
Damage triggers panic & Mindset situation & Damage is information \\
Samira has many buttons & Button-order situation & Q before E; W saved for real danger \\
Low targets bait chase & Chase situation & Low plus reachable, not low plus fog \\
Champion fit affects repetition & Three champion pages & Samira and Cait recordings; Fizz empty until footage exists \\
\bottomrule
\end{tabular}
}
\end{table}

\subsection{The interface itself became part of the research record}

The public page is not a neutral wrapper around the recordings. It is part of the intervention being studied. For that reason, the paper records the actual interface, not only the backend, manifest, or coaching text. The page begins with the shared AO Labs header and quiet paper strip, then presents two compact work surfaces on one continuous page. \emph{Samira tips} accepts pasted, dropped, or selected screenshots and shows the source image with its grounded tips. \emph{Coach game analyses} keeps the separate long-text intake, the coach-message copy control, the current takeaway, two performance charts, and saved-game entries. The earlier champion-card surface used exact user-supplied photographs for Samira, Caitlyn, and Fizz, and those source assets remain in the project archive, but they are no longer part of the default task.

The screenshot path is intentionally independent from the game-entry path. The server saves the exact accepted PNG, JPEG, or WebP before analysis, addresses duplicates by the original SHA-256, and creates a display thumbnail without replacing the source. Each public record carries a stable id, state, transcript, source summary, atomic tip ids, and morning-email eligibility. Pending work survives a restart; unavailable or irrelevant analysis never turns into invented advice. A screenshot may contribute a later practice cue, but it cannot add a performance-rank point, CS point, game fact, or note PDF.

The coach path begins with one button. Pressing \emph{copy coach message} places one complete VOD-review request on the clipboard without changing or submitting the coach-response textarea. The user attaches that text and the video in an external coach conversation, then pastes the completed analysis back into League. The raw response is persisted before deterministic parsing or generated structuring. The parser gives final Alan/Samira scoreboard labels precedence over team, enemy, or interim score lines, and an explicit gameplay-estimated rank remains authoritative for the entry, PDF, trend, and aggregate read. CS@10 stays early-farm evidence rather than a hidden rank input. A structured derivative may organize facts, scoreboard values, rank evidence, timeline decisions, lane, mechanics, fighting, macro, resources, vision, mental and communication patterns, development priorities, drills, and uncertainties. Each derived claim is marked as coach-stated, grounded derivative, or not visible. Failure of that derivative leaves the exact coach response, deterministic facts, current rank read, and PDF intact.

The two chart cards share a row on wide screens instead of centering a narrow fixed-width plot inside a full-width blank card. Their geometry follows the measured container. The tip and note grids use natural-height cards and collapse from three columns to one without stretching short content. The page keeps a 1,280-pixel working maximum, modest gutters, direct labels, and bounded empty regions. Loading does not render empty final cards or blank chart shells. This density is not decorative compression; it keeps the next source action and recent evidence visible in the same reading window.

Each recording row is also a designed reading object, but the current visible contract is no longer a coach paragraph. The row now shows a two-column mistake table only: \emph{Time} and \emph{Mistake}, capped at five rows. The table is intentionally narrow because the user asked to reduce overload: it answers only when the mistake happened and what mistake category occurred. The allowed visible categories are red-light E or red-light commit, second fight after value, and chased into bad space. Rank explanation, broad pressure-mode titles, and internal fallback/model wording are not part of the visible review.

\begin{table}[H]
\centering
\caption{Current interface elements and the paper obligation they create.}
{\small
\setlength{\tabcolsep}{5pt}
\begin{tabular}{p{0.23\linewidth}p{0.33\linewidth}p{0.30\linewidth}}
\toprule
Visible surface & Function in the practice room & Paper record required \\
\midrule
Shared AO Labs header and paper strip & Places the tool inside the AO Labs suite and keeps the PDF reachable without adding navigation clutter & Name the public route, suite context, and paper access pattern \\
Tip-image intake & Preserves public source screenshots and exposes grounded tips without changing game evidence & Record exact-original retention, static-image limits, stable tip ids, retry/delete boundaries, and morning eligibility \\
Coach-message button and textarea & Copies the complete external VOD request and receives the coach's completed response & Record that the website neither uploads nor sends video and that copying cannot submit or change the textarea \\
Saved coach-game entries & Preserve raw coach text, facts, rank evidence, timeline, domains, development plan, uncertainties, and PDF & Record primary-source-first persistence, explicit-rank precedence, final Alan/Samira ownership, extraction failure behavior, coverage, and delete boundary \\
Notes-only charts & Show game-time rank and separate CS@10 evidence from coach entries & Preserve the source boundary and describe container-aware, nonblank chart rendering \\
\texttt{/api/samira/tips} read & Supplies stable eligible tips to the morning queue without an email-time model call & Record image-first rotation, coach and legacy fallback, and post-send raw-MIME history gate \\
Recording archive and manifest & Preserve timestamped review evidence for later analysis & State that recordings are source/archive material, not the default visible page \\
Older champion-card assets & Preserve exact supplied photos and prior fit comparisons & State that other champions are retained as source history unless Alan reopens that surface \\
\bottomrule
\end{tabular}
}
\end{table}

\subsection{Recordings moved behind the active capture surfaces}

The site now organizes the visible page by tip-image intake, coach-analysis intake, and saved coach-game entries rather than by abstract lesson cards, champion choice, recording review, or an archive grid. Situation notes and recording rows still exist because they explain the origin of the rules and preserve evidence, but they are no longer loaded into the default home page. This format was narrowed because the user needed less cognitive load than a giant rule wall, stat audit, champion picker, or paragraph coaching taxonomy.

\begin{table}[H]
\centering
\caption{Example Samira rules retained in the archived situation record.}
{\small
\setlength{\tabcolsep}{6pt}
\begin{tabular}{p{0.42\linewidth}p{0.44\linewidth}}
\toprule
Situation prompt & Combined response \\
\midrule
The fight feels urgent and E wants to become the first button. & Start with Q or an auto so the fight gives information before the dash commits the body. \\
The enemy has CC that is not recognized yet. & Stay outside on Q-only, watch one danger spell, use W or sidestep for visible projectiles, and E only after the scary spell is gone. \\
A low-health enemy runs toward fog or teammates. & Stop at the edge; low only matters when the target is also reachable and cheap to kill. \\
W feels necessary immediately after dashing in. & Hold W for the real spell; it is a parry for visible danger, not armor against being close. \\
R works and the excitement makes staying feel automatic. & Treat R as the reward, then leave unless the next target is already low and reachable. \\
E resets after a takedown. & Read the reset as a new check, not an instruction to dash into tank, tower, fog, or fresh crowd control. \\
One hit lands and feels like proof that the enemy is competent. & Treat damage as information, step back, stabilize the camera, and choose the next action. \\
The player is near a team fight, R is not ready, and waiting there costs HP. & Move to the edge, build style from safety, and enter only when S rank exists and the fight is already committed. \\
The player tunnels on one low target. & Check target, danger, target, danger; stop chasing if the chaser becomes the low target. \\
\bottomrule
\end{tabular}
}
\end{table}

\subsection{Other champions became source history}

The champion comparison material is no longer a visible picker on the home page. It remains source history. Samira remains the active intake target because she provides a high-reward dash-reset-R loop, but she also creates the central panic trap of using E as an anxiety button. Caitlyn remains the comfort-sniper reference for range, farming, and unavailable positioning. Fizz remains the comparison reference because his button shape makes the desired loop clear: mark, enter, stab, dodge, leave.

The older champion-fit material remains useful as source history, but it is not allowed to compete with the current practice page. Both active capture paths are explicitly Samira-scoped, while recording and other-champion feedback stay hidden until Alan reopens those surfaces. This avoids fake specificity and preserves the rule that no footage means no recording feedback.

\subsection{Raw recordings converted carry games into timestamped improvement lanes}

On May 18--25, 2026, the project added a generated recordings manifest from the local League of Legends Highlights folder and the live capture publisher. The current manifest contains 55 recordings across 43 matches and includes short highlight clips plus longer review recordings. Champion detection is bounded to source evidence: older rows remain Samira when sampled frames show Samira in the team list or center the Soul Fighter Samira model and nameplate, while the newest Cait row is source-confirmed from the local client champion id and visible Cait review footage. No champion label is treated as source-confirmed without that kind of evidence.

The recording review does not attempt to become a full coaching dashboard. When reopened, recording rows use a two-column mistake table instead of a narrative paragraph: \emph{Time} gives the clickable game-clock anchor, and \emph{Mistake} gives one short allowed category plus the concrete wrong action and replacement action when visible. Rows are capped at five mistakes so the review answers the immediate question without forcing the user to parse rank explanations, paragraph color semantics, or internal audit language. The performance-rank estimate still exists in the manifest and API source graph, using Riot's current tier structure and queue distinctions \cite{riotRankedTiers2026,riotApexTiers2026}, but it is no longer a home-page surface. Co-op vs. AI rows remain lower-confidence ranked-habit evidence, while ranked/PvP full games remain the stronger transfer evidence. If the detailed model pass fails or returns no usable JSON, the sync builds a conservative table from verified clock frames and hides internal fallback wording from the visible review.

\noindent The May 24 dense-review update keeps that example-review contract but reduces visible repetition and whitespace. Every new current-standard full-game row carries a source-backed victory or defeat label when the League Client match history provides it, uses one sharp mistake title, and starts with one natural key timestamp sentence that names the exact visible state, the wrong click, and the correct next click. Red text is reserved for the leak stated once, green text for one concrete habit to keep, and a distinct warm rose-pink text for one ranked-game \emph{Rep:} drill; the paragraph no longer bolds those spans, so color rather than weight carries the skimmable meaning. The pink drill is now selected from the recording-specific mistake family instead of reused as a universal checklist: death-heavy games drill low-HP exits and no re-entry, death-heavy lane games drill the exact no-dash-under-tower support-and-wave condition, high-kill losses drill cashing out after the first won exchange, grouped base-push wins drill structure--blocker--wave--exit click order, side-farm leaks drill threatened-turret and ally-death checks, and cleaner wins drill the one remaining exit branch while preserving the improvement signal. Neutral text supplies only the timestamped state, consequence, and K/D/A--CS context needed to support those colored claims. At desktop width the recording row uses a balanced text/video split with larger coaching type, so the paragraph and replay are peers rather than a small note beside a dominant video. Rows are rejected when they fall back to generic review-frame language, label lists such as ``mistake category'' and ``correct next click,'' repeated tower--wave--objective--front checklists, a pink drill that does not match the game-specific mistake type, vague ``permanent thing'' drill wording, or abstract branch terms that do not become an if-this-then-that click rule.

\noindent The May 25 forensic-performance update adds a stricter rank and timestamp contract. A row now reports one \emph{Approx performance rank} for that game, never a range, with a one-sentence reason tying together CS/min, K/D/A and death count, fight-entry legality, exit or re-entry quality, and whether pressure became objective, wave, recall, group, or structure value. The key timestamp must separate the play into phases in natural prose: legal or illegal entry, value created or not, state flip, first bad next click, and correct next click. Objective-fight reviews are calibrated so the first entry is not called wrong just because the final result was bad; if dragon or Baron is visible, allies are committed, and enemies are already engaged, the review can mark the entry partly legal and critique the second fight after value instead. The pink drill remains game-specific: lane deaths drill first safe exit and no illegal E, objective fights drill dragon/wave/recall/group after first value, fed losses drill the first won fight into objective/tower/wave/recall, grouped base pushes drill structure/blocker/wave/exit, and jungle-fight exits drill turret, mid wave, reset, or regroup after the first value window. A May 26 micro-feedback correction adds the same specificity requirement to Samira fight rows, not only the early-lane side pass: if the timestamp is a lane trade or mid-fight re-entry, the review must say whether the first entry or damage window is playable, identify the first auto/Q or burst window, name the back-click or reset-spacing point, and state whether the next E or forward click is legal or just a chase. The same branch gate rejects recycled objective, farm, pressure-mode, or macro-only phrasing when the title, timestamp, red leak, and pink rep do not describe the same actual play type.

\noindent The May 27 Samira update narrows future reviews to a single habit until the user changes the standard: green-light discipline for \texttt{E} toward champions. A Samira row must first ask whether all three fight conditions are true: \texttt{W} ready, HP above half, and an ally on screen or close enough to fight with her. If any condition fails, the moment is red light: \texttt{W} gray means no \texttt{E}, below-half HP means no \texttt{E}, and no nearby ally means no \texttt{E}. The review must find the key timestamp where this rule breaks, call the moment red light or green light, state which condition failed, name the wrong input, and translate the correction into a literal input such as Q/auto while backing up, kite, farm, wait, or reset spacing. Macro language about rotations, broad cash-out, or objective conversion is demoted unless it directly follows from a green-light violation. The pink \emph{Rep:} line is therefore no longer allowed to drift into a generic map drill for Samira; it must preserve the next-game cue ``W, HP, ally'' and the operational rule that red-light Samira farms, kites, or waits while green-light Samira can fight.

\noindent The latest May 27 overload-reduction update removes paragraph reviews from the visible recording cards. Future public reviews are mistake tables only, with columns \emph{Time} and \emph{Mistake}, five rows max. The table vocabulary is deliberately small: red-light E or red-light commit, second fight after value, and chased into bad space. The row text must be short and concrete, for example naming which green-light condition failed, that a second fight happened after value, or that a chase moved into tower, brush, river, jungle, or enemy side. Rank explanations, broad ``pressure mode'' labels, and internal fallback/model wording are rejected at the manifest audit layer.

\noindent The later May 25 rank-audit pass recalibrated every full-review row in the 55-recording manifest rather than only the newest sample. The audit rejects Gold-or-higher grades without Gold-level CS/min, death count, queue pressure, and conversion evidence; caps Co-op vs. AI rows below transferable high-rank claims; keeps low-impact or death-heavy games in Iron unless carry-level damage and farm justify a Bronze floor; preserves known Ranked Solo queue metadata from earlier local-client evidence when current logs no longer expose the queue id; and requires each visible rank reason to mention CS/min, deaths or K/D/A, entry or fight context, and exit or value conversion in one sentence. This calibration changed the rank trend by removing inflated Gold/Silver reads and by treating wins, losses, bot games, and high-kill losses as performance evidence rather than outcome labels.

\subsection{The public log made mindset part of the artifact}

The public log starts with source entries for the seed arc import, the practice-room frame, the current Samira law, and later user-supplied Samira proof-game notes. Coach game analyses continue in that persisted note corpus and are exposed for cards, structured entry reads, note PDFs, rank reads, and chart history. Screenshot tips use a different manifest and source directory. Their morning interface returns stable tip ids and source types without letting the automation reinterpret the image at send time. Both paths are intentionally public. The page states that boundary next to screenshot intake, and the coach-response surface remains a public record. Unrelated sensitive information such as credentials, access tokens, or private third-party identifiers remains outside the artifact.

\section{Discussion}

The case suggests that the useful interface for some high-repetition learners is not less information in the abstract, but the right source boundary and the right next action. A screenshot of general Samira advice and a coach's account of one recorded game are both useful, but treating them as one note type would corrupt the performance record. The current design therefore keeps screenshot tips in a public advice library and coach responses in the game-entry corpus. The morning queue can rotate the former and bounded next-game rules from the latter, while the rank and CS charts remain tied only to coach game analyses. The same separation reduces interaction burden: one button prepares the complete request for an external VOD review, one textarea receives the completed response, and one compact image control receives general tips. Dense, natural-height cards and paired charts keep those actions visible without a broad dashboard or large blank containers. This is consistent with cognitive-load theory: reducing simultaneous interpretation and keeping source meaning stable can make practice more tractable \cite{sweller1988}. It is also compatible with mental practice research, because stored decisions can become concrete game cues without making capture feel like another review assignment \cite{driskell1994}.

The work also reframes death and failure. The site does not encourage intentional feeding or fake failure. It uses the rule ``choose the fight; death is allowed; winning is intended.'' That distinction matters. It keeps agency with the learner while preserving the real objective of the game. It also separates practice failure from collapse: a failed engage that tested damage, W timing, or exit timing is a rep; walking in only to die is not.

Several limits are important. First, this is a single-user artifact, not a clinical intervention. The user self-described ADHD/autism traits and supplied a practice narrative, but the site does not diagnose, treat, or generalize to neurodivergent players as a class. Second, the current evidence is artifact-level: source text, distilled rules, local recordings, local League Client match-history reads, a rendered page, and the paper itself. It does not include a controlled comparison or a ranked outcome. Third, the system deliberately avoids public Riot API import in version 1. That keeps the artifact focused on review and self-regulation rather than turning it into a general stat tracker.

\section{Materials and Methods}

\subsection*{Source corpus}

The seed corpus was the user-provided League learning arc pasted into the implementation thread on May 17, 2026. It described the user's own League practice, including early Caitlyn farming goals, Samira drills, bot-ladder practice, camera and sensory adjustments, death exposure, and champion-fit comparisons. Later May 17 and May 18 prompts added Samira crowd-control, W-blocking, poke-lane, bait-recognition, team-fight, bad-team, normal-game, and proof-game refinements. A May 18 local recording folder then supplied 15 files across matches \texttt{NA1-5563247854}, \texttt{NA1-5563301586}, and \texttt{NA1-5563352800}. Local replay files supplied the match-time anchor, League client logs supplied queue id \texttt{880}, Riot's published queue table identifies \texttt{880} as Co-op vs. AI Beginner Bot games,\cite{riotqueues2026} and the local League Client match-history endpoint supplied K/D/A, CS, and game duration when the client was available. The implementation treated these texts, recordings, replay files, local logs, and local client reads as the current source of truth for the public situation notes and recording review.

\subsection*{Interface requirements}

The page was constrained by user instructions gathered during planning and revision: one page only, clean and pretty, the same AO Labs visual family, no poster hero, no generic dashboard, no redundant helper cards, no overexplaining, one quiet paper strip, and no ridiculous white space. Earlier paragraph reviews, champion cards, and recording tables remain source/archive layers rather than the default home page. The current visible interface uses a warm off-white background, dark ink, restrained burgundy and gold accents, a 1,280-pixel maximum content width, and an 8/12/16/24/32-pixel spacing system. The first work surface is a compact public screenshot intake with an adjacent current tip summary only when summary data exists. The second is headed \emph{coach game analyses} and keeps \emph{copy coach message} visible above \emph{paste your coach's completed analysis}. The complete coach paragraph stays hidden unless clipboard access fails. The charts share a row when space allows, and both image and coach-entry cards use natural heights and responsive one-, two-, or three-column grids. The rank path remains source-bounded: explicit rank phrases in the saved coach response win first; parsed final Alan/Samira game facts set the fallback score; CS@10 remains a separate charted fact; ordinary champion names and screenshot tips cannot produce rank points. This interface description is part of manuscript accuracy: when the page changes what Alan sees, submits, interprets, or learns, the paper must change in the same work cycle.

\subsection*{Rule extraction}

Rules were extracted by finding repeated gameplay or mindset triggers in the source corpus and translating each into a one-sentence situation prompt plus one combined response paragraph. The extraction favored recognition over slogan compression: the prompt had to name the exact stressful moment, and the response had to combine action, reason, and practice rep without separate ``do,'' ``avoid,'' or ``drill'' labels. For example, the Samira section retained the underlying rules ``Q first; earn the E,'' ``W is parry, not armor,'' and ``R is reward, not start,'' but the public card now starts from the situation that makes the rule necessary. The death section retained ``choose the fight; death is allowed; winning is intended.'' The bot-practice section retained the practice ladder and the current-mode cue.

\subsection*{Public artifact implementation}

The website is served by a small Node application on Railway. It serves the hand-authored page, recording archive, paper artifacts, and two independent persistent Samira stores. \texttt{/api/samira/notes} accepts the raw coach response; \texttt{/api/samira/notes/:id} reads the complete structured game entry; and \texttt{/api/samira/notes/:id.pdf} renders its document. \texttt{/api/samira/tip-images} accepts one raw image body per request and lists the public screenshot library. Record-specific routes expose the exact original, display thumbnail, detail, bounded retry, and delete operations. \texttt{/api/samira/tips} returns normalized eligible tips with stable ids for the morning queue. The existing \texttt{/api/samira} remains compatible and continues to expose the notes-only aggregate rank, rank trend, CS@10 trend, and compact entry status.

The image store is hash-addressed and separate from the coach-note JSON and coach-entry cache. The browser caps each intake action at five files. The server checks type declarations, magic bytes, decoded dimensions, byte size, animation state, path safety, rate limits, and the 200-record cap before acceptance. A duplicate SHA-256 returns the existing record. One sequential event-driven worker produces a transcript, summary, and source-grounded tips after the original is durable. A record may be pending, ready, or unavailable; analysis retries are bounded, and late worker output is ignored after deletion. The coach-entry derivative is keyed by note id, body hash, and extraction version. Primary-note persistence precedes generated extraction, so provider failure cannot lose the response. Both public write paths are password-free by explicit product choice. The page renders generated text as text rather than markup.

Local recording media and automatic post-game review remain archive sources. The website does not accept a coach video upload and does not contact the coach. The copied request is the complete handoff; the video remains in the user's existing external channel. Public Riot API match history is not imported in this version.

\backmattersection{Acknowledgements}

The source learning arc was supplied by Alan Pham during the implementation conversation. Riot Games was not involved in the design, writing, or deployment of this artifact.

\backmattersection{Funding}

No external funding supported this work.

\backmattersection{Author contributions}

Alan N. Pham supplied the practice corpus, product constraints, and acceptance criteria. The manuscript and site were assembled from those source materials in the AO Labs workspace.

\backmattersection{Competing interests}

The author owns and operates AO Labs. The site is a personal public artifact and is not endorsed by Riot Games.

\backmattersection{Data availability}

The first public data sources are the League practice arc supplied in the implementation thread, exact screenshots intentionally submitted to the public tip library, and coach analyses intentionally pasted into the public game-entry intake at \href{https://league.aolabs.io}{league.aolabs.io}. The home page exposes screenshot thumbnails and grounded tips, coach-game entries, a notes-only game-time rank trend, a separate CS@10 trend, and public entry PDFs. Exact screenshot originals remain available through their source routes until publicly deleted. The API also exposes stable morning-tip ids and source types, deterministic coach-game facts, structured entry derivatives with uncertainty labels, generated recording manifest data, poster frames, local queue metadata, local client K/D/A and CS when available, per-recording mistake tables, compact post-game queue status, and timestamp anchors for archived videos. Screenshot tips are excluded from rank and CS evidence. The external coach video itself is not received, stored, or published by this website. Public Riot API match history is not imported in this version.

\backmattersection{Code availability}

The implementation is intended for the \texttt{nalalalan/league-app} repository, with the public site, PDF, manuscript source, and log API served at \href{https://league.aolabs.io}{league.aolabs.io}.

\backmattersection{Additional information}

league.aolabs.io is not endorsed by Riot Games and does not reflect the views or opinions of Riot Games or anyone officially involved in producing or managing Riot Games properties. Riot Games, League of Legends, and associated properties are trademarks or registered trademarks of Riot Games, Inc.

\printbibliography[title={References}]

\end{document}

